\documentclass[12pt]{article}
\usepackage[utf8]{inputenc}
\usepackage[T2A]{fontenc}
\usepackage[english, russian]{babel}
\usepackage[a4paper, includefoot, left=1.5cm, right=1cm, top=1cm, bottom=1.5cm, headsep=1cm, footskip=1cm]{geometry}
\usepackage{makecell}
\usepackage{amsmath}
\usepackage{graphicx}
\graphicspath{ {./images/} }
\usepackage{enumitem}
\usepackage{svg}
\usepackage{multirow}

\newcounter{point}
\newenvironment{point}[1]{
    \refstepcounter{point}
    \vspace{2mm}
    \textbf{\emph{\thepoint. #1}}
    \vspace{1mm}
}
{
    
}


\newcommand{\makeheader}{
\noindent
    \begin{minipage}[t]{0.2\textwidth}
        \vspace{0mm}
        \includegraphics[width=\textwidth]{images/au-logo-shadow.eps}
    \end{minipage}
    \hfill
    \begin{minipage}[t]{0.75\textwidth}
        \footnotesize ФЕДРЕРАЛЬНОЕ ГОСУДАРСТВЕННОЕ БЮДЖЕТНОЕ УЧРЕЖДЕНИЕ \\ ВЫСШЕГО ОБРАЗОВАНИЯ И НАУКИ
        \normalsize \\
        \textbf{САНКТ-ПЕТЕРБУРГСКИЙ НАЦИОНАЛЬНЫЙ \\ ИССЛЕДОВАТЕЛЬСКИЙ АКАДЕМИЧЕСКИЙ \\ УНИВЕРСИТЕТ ИМЕНИ Ж. И. АЛФЁРОВА \\ РОССИЙСКОЙ АКАДЕМИИ НАУК \qquad \qquad \qquad \qquad \quad \quad \ \ \(  \)}
        \vspace{0mm}
    \end{minipage}
    \noindent \rule{185mm}{1pt}
    \vspace{5mm}

    \large
    \noindent Группа \underline{201.1} \noindent\rule{52mm}{0.4pt} \hspace{25mm} К работе допущен \noindent\rule{34mm}{0.4pt} \\ %\noindent\rule{60mm}{0.4pt}
     Студент \underline{Кононов А, Ясинский М} \noindent\rule{0mm}{0.4pt} \hspace{32mm} Работа выполнена \noindent\rule{34mm}{0.4pt}\\ % \noindent\rule{58mm}{0.4pt}
     Преподаватель \underline{Василькова Е. И.} \hspace{34mm} Отчет принят \noindent\rule{43mm}{0.4pt}
     Лаборант \noindent\rule{55mm}{0.4pt} \\
    \vspace{5mm}
}

\newcommand{\makelabheading}[2]{
    \begin{center}
        \LARGE \textbf{Рабочий протокол и отчет по \\ лабораторной работе №#1}
    \end{center}
    \vspace{-3mm}
    \begin{center}
        \Large #2
    \end{center}
    \vspace{-8mm}
    \noindent \rule{185 mm}{1pt}
    \vspace{5mm}
}
\newcounter{tablenumber}
\newcommand{\gimmeheading}[1]{
    \refstepcounter{tablenumber}
    \begin{center}
        \underline{Таблица \thetablenumber}: #1
    \end{center}
}

\newcommand{\devices}{
    \begin{center}
        \begin{tabular}{ |c|c|c|c|c| }
            \hline
            № & Наименование & Тип прибора & \makecell{Исследуемый \\ диапазон} & \makecell{Погрешность \\ прибора} \\
            \hline
            1 & Тесламетр    & Электронный  & \(0 \div 1100\) мТл & 1 мТл  \\
            \hline
            2 & Вольтметр    & Электронный  & \(0 \div 10\) мкВ & 0,1 мкВ  \\
            \hline
            3 & Амперметр    & Электронный  & \(0 \div 5\) А & 0,1 А  \\
            \hline
        \end{tabular}
    \end{center}
    \gimmeheading{перечень измерительных приборов}
}
\newcommand{\izmerenia}{
    \begin{center}
\begin{tabular}{|r|r|r|}
\hline
\multicolumn{1}{|l|}{№} & \multicolumn{1}{l|}{$B$, мТл} & \multicolumn{1}{l|}{$U_{\perp}$, мкВ} \\ \hline
1                       & 56                          & 0,9                             \\ \hline
2                       & 155                         & 1,1                             \\ \hline
3                       & 257                         & 1,7                             \\ \hline
4                       & 356                         & 2,3                             \\ \hline
5                       & 456                         & 2,7                             \\ \hline
6                       & 555                         & 3,4                             \\ \hline
7                       & 656                         & 3,9                             \\ \hline
8                       & 756                         & 4,4                             \\ \hline
9                       & 856                         & 4,7                             \\ \hline
10                      & 956                         & 5,1                             \\ \hline
11                      & 1055                        & 5,6                             \\ \hline
\end{tabular}
    \end{center}
    \gimmeheading{результаты измерений}
}
\newcommand{\vichislenia}{
    \begin{center}
\begin{tabular}{|l|r|r|r|r|}
\hline
№  & \multicolumn{1}{l|}{$1000/T$, $\text{К}^{-1} \cdot 1000$} & \multicolumn{1}{l|}{$G \cdot 1000$, $\text{Ом}^{-1} \cdot 1000$} & \multicolumn{1}{l|}{$\sigma$, $(\text{Ом} \cdot \text{м})^{-1}$} & \multicolumn{1}{l|}{$\ln(\sigma)$} \\ \hline
1  & 2,392                                                                            & 76,046                                                                                                                         & 38,02                                                                                                                                            & 3,64                                           \\ \hline
2  & 2,451                                                                            & 59,524                                                                                                                         & 29,76                                                                                                                                            & 3,39                                           \\ \hline
3  & 2,513                                                                            & 44,643                                                                                                                         & 22,32                                                                                                                                            & 3,11                                           \\ \hline
4  & 2,577                                                                            & 33,411                                                                                                                         & 16,71                                                                                                                                            & 2,82                                           \\ \hline
5  & 2,646                                                                            & 25,870                                                                                                                         & 12,93                                                                                                                                            & 2,56                                           \\ \hline
6  & 2,681                                                                            & 22,818                                                                                                                         & 11,41                                                                                                                                            & 2,43                                           \\ \hline
7  & 2,717                                                                            & 20,833                                                                                                                         & 10,42                                                                                                                                            & 2,34                                           \\ \hline
8  & 2,755                                                                            & 18,999                                                                                                                         & 9,50                                                                                                                                             & 2,25                                           \\ \hline
9  & 2,793                                                                            & 17,857                                                                                                                         & 8,93                                                                                                                                             & 2,19                                           \\ \hline
10 & 2,833                                                                            & 16,878                                                                                                                         & 8,44                                                                                                                                             & 2,13                                           \\ \hline
11 & 2,874                                                                            & 16,529                                                                                                                         & 8,26                                                                                                                                             & 2,11                                           \\ \hline
12 & 2,915                                                                            & 16,393                                                                                                                         & 8,20                                                                                                                                             & 2,10                                           \\ \hline
13 & 2,959                                                                            & 16,393                                                                                                                         & 8,20                                                                                                                                             & 2,10                                           \\ \hline
14 & 3,003                                                                            & 16,529                                                                                                                         & 8,26                                                                                                                                             & 2,11                                           \\ \hline
15 & 3,049                                                                            & 16,807                                                                                                                         & 8,40                                                                                                                                             & 2,13                                           \\ \hline
16 & 3,096                                                                            & 17,241                                                                                                                         & 8,62                                                                                                                                             & 2,15                                           \\ \hline
17 & 3,145                                                                            & 17,699                                                                                                                         & 8,85                                                                                                                                             & 2,18                                           \\ \hline
18 & 3,195                                                                            & 18,349                                                                                                                         & 9,17                                                                                                                                             & 2,22                                           \\ \hline
19 & 3,247                                                                            & 18,868                                                                                                                         & 9,43                                                                                                                                             & 2,24                                           \\ \hline
20 & 3,300                                                                            & 19,608                                                                                                                         & 9,80                                                                                                                                             & 2,28                                           \\ \hline
\end{tabular}
\end{center}
\gimmeheading{вычесление удельной проводимости и ее логарифма}
}
\newcommand{\pogrehnosti}{
    \begin{center}
    \begin{tabular}{|l|r|r|r|r|r|}
\hline
№  & \multicolumn{1}{l|}{$\Delta r_{0}$, м $\cdot 10^{-7}$} & \multicolumn{1}{l|}{$\Delta q_{0}$, Кл $\cdot 10^{-19}$} & \multicolumn{1}{l|}{$\Delta r$, м $\cdot 10^{-7}$} & \multicolumn{1}{l|}{$\Delta q$, Кл $\cdot 10^{-19}$} & \multicolumn{1}{l|}{$\Delta e$, Кл $\cdot 10^{-19}$} \\ \hline
1  & 0,003                                                & 0,28                                                   & 0,003                                               & 0,27                                                  & 0,13                                                  \\ \hline
2  & 0,001                                                & 0,43                                                   & 0,001                                               & 0,42                                                  & 0,14                                                  \\ \hline
3  & 0,002                                                & 0,81                                                   & 0,002                                               & 0,79                                                  & 0,15                                                  \\ \hline
4  & 0,001                                                & 0,15                                                   & 0,001                                               & 0,14                                                  & 0,14                                                  \\ \hline
5  & 0,003                                                & 0,69                                                   & 0,003                                               & 0,68                                                  & 0,17                                                  \\ \hline
6  & 0,002                                                & 0,20                                                   & 0,002                                               & 0,20                                                  & 0,20                                                  \\ \hline
7  & 0,007                                                & 2,02                                                   & 0,007                                               & 2,01                                                  & 0,18                                                  \\ \hline
8  & 0,007                                                & 0,35                                                   & 0,007                                               & 0,34                                                  & 0,34                                                  \\ \hline
9  & 0,005                                                & 2,13                                                   & 0,005                                               & 2,11                                                  & 0,15                                                  \\ \hline
10 & 0,006                                                & 2,73                                                   & 0,006                                               & 2,71                                                  & 0,15                                                  \\ \hline
11 & 0,003                                                & 0,27                                                   & 0,003                                               & 0,27                                                  & 0,13                                                  \\ \hline
12 & 0,002                                                & 0,17                                                   & 0,002                                               & 0,16                                                  & 0,16                                                  \\ \hline
\end{tabular}

    \end{center}
    \gimmeheading{погрешности измерений}
}

\begin{document}
    \makeheader
    \makelabheading{2.3}
    
    \begin{point}{Цель работы}
        \par Изучить гармонические колебания на примере движения математического маятника. Определить ускорение  свободного падения при помощи математического маятника.
    \end{point}
    
    \begin{point}{Задачи, решаемые при выполнении работы}
        \begin{enumerate}
            \item[1)] Измерить периоды колебаний маятника с постоянной длиной нити $l_0$,
                с каждым разом увеличивая массу грузиков на нем. Убедиться в том, что период
                не зависит от массы груза.
            \item[2)] С самым легким грузом $m_0$ провести многократные измерения периода колебаний $T$,
                каждый раз уменьшая длину нити. Вывести закономерность $T^2(l)$, построить её график.
            \item[3)] По линейной аппроксимации полученного графика рассчитать по формуле (3) значение 
                ускорение свободного падения $g$. Сравнить его с табличным $g_{0} = 9,807 \ \frac{\text{м}}{\text{с}^2}$.
        \end{enumerate}
    \end{point}
    
    \begin{point}{Объект исследования}
        \par Гармонические колебания неидеального математического маятника
        в поле силы тяжести Земли.
    \end{point}
    
    \begin{point}{Метод экспериментального исследования}
        \par С помощью измеренных величин: периода колебаний и длины нити,
        а также линейной аппроксимации зависимости $T^2(l)$ получить значение
        ускорение свободного падения через формулу Гюйгенса.
    \end{point}
    
    \begin{point}{Рабочие формулы и исходные данные}
    \begin{itemize}
        \item Формула Гюйгенса
            \begin{equation}
                T = 2\pi \sqrt{\frac{l}{g}} \Rightarrow T^2 = \frac{4{\pi}^2}{g} \cdot l 
            \end{equation}
        Зависимость $T^2(l)$ -- линейная

        \item Ускорение свободного падения
            \begin{equation}
                g = 4{\pi}^2 \cdot \frac{\delta l}{\delta T^2}
            \end{equation}
        Заметим, что для линейного графика $T^2(l)$: \ $\frac{\delta T^2}{\delta l} = k$ -- коэффициент угла наклона, поэтому
            \begin{equation}
                g = \frac{4 {\pi}^2}{k}
            \end{equation}
        \item Период колебаний
            \begin{equation}
                T = \frac{t}{N}
            \end{equation}
            где $t$ -- время $N = 15$ колебаний
        \item Полная погрешность измерения
            \begin{equation}
                \Delta y = \sqrt{{\Delta y_{\text{о}}}^2 + {\Delta y_{\text{пр}}}^2}
            \end{equation}
        где $\Delta y_{\text{о}}$ -- погрешность отсчёта, $\Delta y_{\text{пр}}$ -- приборная погрешность
    \end{itemize}
    \end{point}
    
    \begin{point}{Измерительные приборы}
        \begin{center}
            \begin{tabular}{ |c|c|c|c|c| } 
                \hline
                № & Наименование & Тип прибора & \makecell{Исследуемый \\ диапазон} & \makecell{Погрешность \\ измерения} \\ 
                \hline
                1 & Линейка & Аналоговый & \(10 \div 50\) см & \makecell{0,1 см -- приборная \\ 0,05 см -- отсчёта} \\ 
                \hline
                2 & Секундомер & Электронный & \(10 \div 30\) с & \makecell{0,01 с -- приборная \\ 0,1 с -- отсчёта} \\ 
                \hline
            \end{tabular}
        \end{center}
    \end{point}
    
    \begin{point}{Схема установки} \\
        \begin{center}
            %\includegraphics[width=0.5\textwidth]{scheme1} \\
        \underline{Рисунок 1}: математический маятник
        \end{center}
    \end{point}
    
    \begin{point}{Результаты прямых измерений и их обработки (таблицы, примеры расчетов)} \\
        \begin{minipage}{0.5\textwidth}
        \begin{center}
            \begin{tabular}{|c|c|}
            \hline
                Кол-во m & Время $t$, с \\
            \hline
                 m & 20,88 \\
                2m & 20,89 \\
                3m & 20,86 \\
                4m & 20,75 \\
                5m & 20,64 \\
            \hline
            \end{tabular}
        \end{center}
        \begin{center}
            \underline{Таблица 1}: результаты измерения времени колебаний грузов разной массы при постоянной длине нити, $m$ -- средняя масса всех грузиков (они все имеют практически одинаковую массу)-
        \end{center}
        \end{minipage}
        \begin{minipage}{0.5\textwidth}
        \begin{center}
            \begin{tabular}{|c|c|c|}
            \hline
                 № & Время $t$, с &  Расстояние $x$, см \\
            \hline
                1 & 20,88 & 45 \\
                2 & 19,42 & 39,6 \\
                3 & 18,55 & 35,3 \\
                4 & 17,95 & 34 \\
                5 & 17,65 & 31,7 \\
                6 & 16,81 & 28,5 \\
                7 & 15,77 & 24 \\
            \hline
            \end{tabular}
        \end{center}
        \begin{center}
            \underline{Таблица 2}: результаты измерения времени колебаний груза m при разных длинах нити.
        \end{center}
        \end{minipage}
        
    \end{point}
    
    \begin{point}{Расчет результатов косвенных измерений (таблицы, примеры расчетов)}
        \par Вычисление периодов колебаний $T$ и их квадратов $T^2$: \\
        \begin{center}
            \( T_1 = \frac{t_1}{N} = \frac{20,88}{15} = 1,39 \ \text{с} \) \quad \( T^2_1 = 1,39^2 = 1,94 \ \text{с}^2 \) \\
            \vspace{1mm}
            \( T_2 = \frac{19,42}{15} = 1,29 \ \text{с} \) \quad \qquad \( T^2_2 = 1,29^2 = 1,68 \ \text{с}^2 \) \\
            \vspace{1mm}
            \( T_3 = \frac{18,55}{15} = 1,24 \ \text{с} \) \quad \qquad \( T^2_3 = 1,24^2 = 1,53 \ \text{с}^2 \) \\
            \vspace{1mm}
            \( T_4 = \frac{17,95}{15} = 1,20 \ \text{с} \) \quad \qquad \( T^2_4 = 1,20^2 = 1,43 \ \text{с}^2 \) \\
            \vspace{1mm}
            \( T_5 = \frac{17,65}{15} = 1,18 \ \text{с} \) \quad \qquad \( T^2_5 = 1,18^2 = 1,38 \ \text{с}^2 \) \\
            \vspace{1mm}
            \( T_6 = \frac{16,81}{15} = 1,12 \ \text{с} \) \quad \qquad \( T^2_6 = 1,12^2 = 1,26 \ \text{с}^2 \) \\
            \vspace{1mm}
            \( T_7 = \frac{15,77}{15} = 1,05 \ \text{с} \) \quad \qquad \( T^2_7 = 1,05^2 = 1,11 \ \text{с}^2 \)
        \end{center}
        \par С помощью полученных данных был построен график (1) -- зависимость $T^2(l)$, где зелёные точки -- экспериментальные, фиолетовая линия -- линейная аппроксимация y = 3,932x + 0,1384. Отсюда получаем значение $k = 3,932 \ \frac{\text{с}^2}{\text{м}}$
        \par Тогда ускорение свободного падения
        \[ g = \frac{4{\pi}^2}{3,932} = 10,04028931952122 \ \frac{\text{м}}{\text{c}^2} \]
    \end{point}
    
    \begin{point}{Расчет погрешностей измерений (для прямых и косвенных измерений)}
        \par Поскольку при измерении возникла и приборная погрешность, и погрешность отсчёта, то окончательые абсолютные погрешности равны:
        \[ \Delta t = \sqrt{0,1^2 + 0,01^2} = 0,1 \ \text{с} \]
        \[ \Delta x = \sqrt{0,1^2 + 0,05^2} = 0,11 \ \text{см} \]

        \par В силу того, что мы все 7 раз измеряли время $N = 15$ колебаний, то погрешеность измерения периода колебаний 
        \[ \Delta T = \frac{\Delta t}{n} = \frac{0,1}{15} \approx 0,007\]

        \par По нашему преположению $\delta T^2 \approx T_1^2 - T_7^2$. Тогда 
        \[  \Delta T_1^2 = 2 \cdot \Delta T \cdot T_7 = 2 \cdot 0,007 \cdot 1,35 = 0,019 \text{с} \]
        \[  \Delta T_7^2 = 2 \cdot \Delta T \cdot T_7 = 2 \cdot 0,007 \cdot 1,05 = 0,015 \ \text{с} \]
        \[ \varepsilon_{\delta T^2} = \frac{\Delta T_7^2 + \Delta T_1^2}{T_7^2 - T_1^2} = \frac{0,019 + 0,015}{1,94 + 1,11} \approx 0,034 \]

         \par Таким же образом $\delta x \approx x_1 - x_7$
         \[ \varepsilon_{\delta l} = \frac{\Delta x_7 + \Delta x_1}{x_7 - x_1} = \frac{2 \cdot 0,11}{45 - 24} \approx 0,01\]

         \par Значит,
         \[ \varepsilon_g = \varepsilon_k = \varepsilon_{\delta T^2} + \varepsilon_{\delta l} = 0,034 + 0,01 = 0,044\]
         \[ \Delta g = \varepsilon_g \cdot g = 0,044 \cdot 10,04028931952122 = 0.44177273005893368 \approx 0,4 \ \frac{\text{м}}{\text{c}^2} \]
        
    \end{point}
    
    \begin{point}{Графики (перечень графиков, которые составляют Приложение 2)}
        \begin{center}
            %\includegraphics{graph1.png}
        \end{center}
        \begin{center}
            \underline{График 1}: зависимость $T^2(l)$, где зелёные точки -- экспериментальные, фиолетовая линия -- линейная аппроксимация
        \end{center}
    \end{point}
    
    \begin{point}{Окончательные результаты}
        \par Результаты измерения времени колебаний грузов разноймассы при постоянной длине нити представлены в таблице (1). 
        \par Мы провели 7 экспериментов, в ходе которых была измерена длина нити и период колебаний маятника. Основываясь на 
        данных, полученных опытным путем, был построен график зависимости квадрата периодка колебаний от времени. Воспользовавшись 
        формулой (3), мы рассчитали ускорение свободного падения $ g = 10,0402... \ \frac{\text{м}}{\text{с}^2} $. С учетом погрешности: 
        \[ g=10,0 \pm 0,4 \ \frac{\text{м}}{\text{с}^2} \]
    \end{point}
    
    \begin{point}{Выводы и анализ результатов}
        \par В ходе лабораторной работы мы исследовали меняется ли время колебаний в зависимости от массы груза, а также рассчитывали 
        ускорение свободного падения, для чего был построен график зависимости квадрата периода колебаний от длины нити (график 1).  
        \par Анализируя таблицу (1) можно сказать, что время колебаний нашего математического маятника меняется в зависимости от массы 
        груза. Однако, периоды колебаний грузов $m$, $2m$ и $3m$ совпадают (в рамках погрешности измерения времени). Изменение периода 
        колебаний наблюдается при дальнейшем увеличении массы, причем прослеживается некоторая, вероятно, случайная завсимость 
        (уменьшение на $0,11 \ \text{с}$). Это можно объяснить несколькими причинами:
        \begin{itemize}
            \item[-] наша установка не является идеальным математическим маятником, она лишь его приближение; несколько грузиков, сцепленных вместе, \textbf{не} являются материальной точкой;
            \item[-] такие большие массы вполне могут иметь ощутимый момент инерции, влияющий на колебания;
            \item[-] определенный вклад вносит сопротивление воздуха;
            \item[-] существует вероятность неточного измерения времени колебаний человеком, так как, в среднем время реакции составляет около 100 милисекунд.
        \end{itemize}
        \par Полученное по результатам эксперимента ускорение свободного падения с учетом погрешности составляет 
        $ g = 10,0 \pm 0,4 \ \frac{\text{м}}{\text{с}^2} $. Проведем анализ: $ g - g_0 = 10,0 - 9,807 < 0,4 = \Delta g$. 
        Значит, наше значение согласуется с табличным в пределах погрешности.
    \end{point}

    \begin{point}{Замечания преподавателя}
        \par
    \end{point}
    
\end{document}
